\documentclass[a4paper]{article}
%\documentclass{report} % Can be article, report, book, anything you want!

\usepackage{hyperref}
\usepackage{vub} % this will automatically set to A4 paper. We're not in the US.
% if, for some obscure reason, you need another format, please file an issue (see further).

% The `vub` package has many options.
% Please read the documentation at https://gitlab.com/rubdos/texlive-vub if you want to learn about them.

% Feel free to mix and match with other packages.

\setlength{\parskip}{1em}
\title{Pharo project}
\subtitle{Metaprogramming and Reflection}
\faculty{Sciences and Bio-Engineering Sciences} % Note: without the word "Faculty"!
\author{Gérard Lichtert\\
        0557513\\
    \href{mailto:gerard.lichtert@vub.be}{gerard.lichtert@vub.be}
}
\date{\today}

\begin{document}
\maketitle
\tableofcontents
\newpage
\raggedright

\section{Overview}
This project consists of three parts.
\begin{itemize}
	\item Protected Visibility Modifier
	\item Object-Centric Debugging
	\item Multiple exception handlers
\end{itemize}
Each part has it's own .st file located in the /src directory.
Do note, that I haven't overriden Pharo's default behavior. For instance, if you
would like to use the protected Visbility Modifier, you will have to add the
plugins to the compiler prior to testing it out. To make testing easier, there
are already tests present that should cover the minimum implementation
requirements.
\section{Protected Visibility Modifier}
\subsection{Requirement 1}
The first requirement was to extend the compiler to support compilation of
$<$protected$>$ methods, using double registration of public methods and selector
mangling of self sends.\par
To achieve this, first we created a plugin for the compiler that mangles the
selector of methods that have the $<$protected$>$ pragma.
We also need double registration, to achieve this, I traced the method install to the point
where methods are installed in a class. This was in ClassDescription, which is
moved to the package since the method to install methods was modified.
\par
A method is installed with, among
others, a selector and a compiledMethod, to which we can ask if it has a
$<$protected$>$ pragma. If it does, we know it will already have a mangled
selector. In this case it should be installed normally.
However, if it does not have the $<$protected$>$ pragma, we know it is a public
method and therefore has to be installed twice, once with the normal selector
and once with the mangled selector. Since a selector is required to be passed
along prior to installation, we install it once normally and install it once but
with the mangled selector. Essentially we add another message send but with the
mangled selector.
\par
To mangle self and super send, we create another plugin, just like we already
did for the selector manging of protected methods. However, this time we will
mangle the messages sent to self and super receivers.
\subsection{Requirement 2}
Another requirement is that the modifier should be backwards compatible, meaning
that sometimes objects will already exist without the mangled selector, or
taking the $<$protected$>$ modifier into account. To overcome this, we need to
be extra careful when mangling self and super sends. More concretely, the plugin
will now check if the class (or superclass, depending on the receiver) has the
mangled method. If it does, it will mangle it, otherwise it will leave it as is.
\section{Object-Centric Debugging}
\subsection{Requirement}
\section{Multiple exception handlers}
\subsection{Requirement 1}
\subsection{Requirement 2}
\subsection{Requirement 3}
\subsection{Requirement 4}

\end{document}
